\chapter{Estado del arte y decisiones arquitectónicas}
\label{cap:estado_arte}

El presente capítulo desarrolla el marco tecnológico necesario para
justificar la arquitectura de la plataforma construida en este Trabajo Fin
de Máster. Se recorre cada capa en el orden en que se despliega:
almacenamiento y formato de tabla, sustrato de cómputo, procesamiento
distribuido, orquestación de flujos, y consulta y visualización. Para cada
una se presentan las alternativas reales del ecosistema y se justifica la
opción adoptada frente a ellas. El capítulo cierra con una síntesis que
reúne todas las decisiones en una tabla y sirve de puente hacia el
Capítulo~\ref{cap:infraestructura}, donde esta arquitectura se despliega
efectivamente.

\section{De los sistemas de almacenamiento analítico al lakehouse}
\label{sec:estado_arte_lakehouse}

Un sistema que integra currículos CVN y perfiles de investigador ORCID
necesita almacenar datos con dos exigencias que tradicionalmente se han
resuelto con arquitecturas distintas: la escala y variedad propias de un
\textit{data lake}, y las garantías transaccionales de un \textit{data
warehouse}.

Un almacén de datos clásico, o \textit{data warehouse}, define el esquema
antes de cargar los datos, una disciplina conocida como
\textit{schema-on-write}. Esto aporta consistencia y consultas analíticas
fiables, justo lo que necesita este proyecto para publicar indicadores. Su
inconveniente aparece antes: ingerir una fuente nueva exige transformarla
primero para que encaje en ese esquema, un proceso costoso cuando CVN y
ORCID llegan en formatos semiestructurados que no coinciden entre sí.

Un lago de datos clásico, o \textit{data lake}, resuelve el problema
contrario. Almacena cualquier fuente tal cual llega y solo interpreta su
esquema al leerla, una práctica conocida como \textit{schema-on-read}. Sin
metadatos transaccionales, sin embargo, corre el riesgo documentado del
\textit{data swamp}: lecturas inconsistentes durante escrituras
concurrentes, y ninguna forma de auditar cómo estaban los datos en un
momento anterior.

El \textit{lakehouse} combina ambas exigencias añadiendo una capa de
metadatos transaccionales, un formato de tabla abierto, sobre
almacenamiento de objetos barato. Esa capa aporta transacciones ACID,
evolución de esquema controlada y \textit{time travel}, sin mover los
datos a un motor relacional cerrado. Armbrust et al.~\cite{armbrust_lakehouse_2021}
formulan esta convergencia como respuesta directa a las limitaciones de
los dos modelos anteriores. La presentan como la arquitectura de
referencia cuando conviven fuentes documentales y de volumen masivo,
exactamente la situación de este proyecto.

La arquitectura adoptada es, por tanto, un \textit{lakehouse} organizado
en tres capas: \textit{bronze}, \textit{silver} y \textit{gold}. La capa
\textit{bronze} conserva la procedencia de cada fuente. La capa
\textit{silver} aplica resolución de identidad y control de calidad. La
capa \textit{gold} deja agregados listos para consumo. Este es el patrón
de fusión multi-fuente que exige el problema planteado en el
Capítulo~\ref{cap:introduccion}. Mantener un único sistema para ambas
exigencias, en lugar de un lago para ingesta y un almacén separado para
publicación, evita además duplicar la superficie operativa del proyecto.

\section{Formato de tabla: Apache Iceberg frente a sus alternativas}
\label{sec:estado_arte_iceberg}

El formato de tabla es la pieza que hace posible un \textit{lakehouse}:
una capa de metadatos sobre ficheros en almacenamiento de objetos que los
convierte en algo que se comporta como una tabla, con transacciones y
evolución de esquema. Tres formatos tienen adopción real en el ecosistema
actual.

Apache Iceberg~\cite{apache_iceberg_docs} nació en Netflix para superar
los límites de las tablas Hive a escala de petabytes, y hoy es un proyecto
de la Apache Software Foundation con gobierno multi-vendedor. Aporta
evolución de esquema y de partición sin reescribir los datos existentes, y
aislamiento de \textit{snapshots} que permite leer una tabla de forma
consistente mientras otro proceso escribe en ella. Cuenta además con
soporte amplio en varios motores de consulta, entre ellos Spark, Trino y
Flink, lo que significa que el mismo catálogo podría consultarse en el
futuro desde un motor distinto sin cambiar el formato de tabla.

Delta Lake~\cite{delta_lake_docs} nació en Databricks y fue donado a la
Linux Foundation. Su integración con el ecosistema Spark de Databricks es
muy pulida, pero sus funciones más avanzadas alcanzan su rendimiento
completo sobre el \textit{runtime} propietario de Databricks. Fuera de ese
entorno, su soporte multi-motor ha sido históricamente más limitado.
Apache Hudi~\cite{apache_hudi_docs} nació en Uber y optimiza
actualizaciones incrementales de baja latencia sobre flujos continuos,
mediante \textit{upserts} y captura de cambios de datos, o CDC. Esta
ventaja no aplica a este proyecto: es un \textit{pipeline} por lotes que
reconstruye \textit{silver} y \textit{gold} por completo en cada
ejecución.

La Tabla~\ref{tab:tfm_formatos_tabla} resume las tres opciones.

\begin{table}[H]
  \centering
  \small
  \renewcommand{\arraystretch}{1.2}
  \begin{tabular}{p{0.16\textwidth} p{0.4\textwidth} p{0.36\textwidth}}
    \hline
    \textbf{Formato} & \textbf{Origen y gobierno} & \textbf{Ajuste con este proyecto} \\
    \hline
    Apache Iceberg &
      Netflix, hoy proyecto Apache, gobierno multi-vendedor &
      Sin dependencia de un proveedor, consultable desde otros motores sin
      cambiar el formato de tabla \\
    Delta Lake &
      Databricks, donado a la Linux Foundation &
      Funciones avanzadas afinadas para el \textit{runtime} de Databricks,
      menor soporte multi-motor fuera de él \\
    Apache Hudi &
      Uber, orientado a actualizaciones incrementales de baja latencia &
      Optimizado para CDC y \textit{streaming}, este proyecto es un
      \textit{pipeline} por lotes \\
    \hline
  \end{tabular}
  \caption{Comparación de formatos de tabla abierta para arquitecturas \textit{lakehouse}.}
  \label{tab:tfm_formatos_tabla}
\end{table}

La opción adoptada es Apache Iceberg. Al no depender de un
\textit{runtime} comercial, las tablas creadas aquí seguirían siendo
consultables sin cambios si, como trabajo futuro, la plataforma migrase a
un motor de consulta interactivo como Trino, tratado en la
Sección~\ref{sec:estado_arte_bi}, o a un clúster gestionado en la nube.
Delta Lake queda como alternativa razonable, pero condicionada a un
entorno Databricks que este proyecto no usa. Hudi se descarta por no
encajar con un \textit{pipeline} por lotes.

\subsection{Catálogo de Iceberg}
\label{sec:estado_arte_catalogo}

Iceberg necesita además un catálogo que resuelva el nombre de una tabla a
la ubicación de su metadato raíz. Las opciones habituales son:

\begin{itemize}
  \item \textbf{Hive Metastore~\cite{apache_hive_metastore_docs}}:
  servicio centralizado que almacena la estructura de tablas y
  particiones. Exige desplegar y operar un servicio adicional únicamente
  para esa resolución de nombres.
  \item \textbf{Catálogo REST~\cite{apache_iceberg_docs}}: la vía que el
  propio proyecto Iceberg promueve como estándar de interoperabilidad
  entre motores. Resuelve el mismo problema mediante una API HTTP común,
  pero exige igualmente un servicio adicional.
  \item \textbf{Project Nessie~\cite{project_nessie_docs}}: catálogo
  versionado al estilo Git, con ramas y \textit{commits} sobre el catálogo
  completo. Muy atractivo para múltiples escritores que necesitan aislar
  sus cambios, pero comparte la misma exigencia operativa.
  \item \textbf{Catálogo Hadoop}, basado en rutas: resuelve el nombre de
  tabla directamente contra una convención de directorios en el propio
  almacenamiento de objetos, sin ningún servicio adicional.
\end{itemize}

La opción adoptada es el catálogo Hadoop. Con un único escritor por capa y
sin necesidad de versionar el catálogo a nivel de repositorio, el coste
operativo de un servicio adicional no se compensa con ninguna capacidad
que el proyecto vaya a usar. El versionado de Nessie y la
interoperabilidad de un catálogo REST quedan como trabajo futuro. Ningún
escenario del proyecto las necesita todavía, no por ninguna carencia
técnica de esas dos opciones.

\section{Almacenamiento de objetos: MinIO frente a sus alternativas}
\label{sec:estado_arte_minio}

HDFS~\cite{apache_hadoop_docs}, la opción clásica de Big Data, acopla
almacenamiento y cómputo en los mismos nodos, lo que contradice el propio
patrón \textit{lakehouse} y añade una pieza operativa adicional sin
aportar ventaja alguna aquí. El almacenamiento nativo de un proveedor de
nube pública es la opción de producción real, pero es incompatible con un
despliegue autoalojado y de coste cero. MinIO~\cite{minio_github}, un
servidor de objetos autoalojado compatible con la API de Amazon S3,
resuelve ambos problemas a la vez: el mismo protocolo, \texttt{s3a://},
que usarían los trabajos de procesamiento contra un proveedor real
funciona sin cambios contra MinIO.

La opción adoptada es MinIO, precisamente por esa compatibilidad de
protocolo. El código de acceso a almacenamiento escrito contra MinIO no
necesitaría cambios para apuntar a un proveedor de nube real.

\section{Sustrato de orquestación de contenedores y entorno de despliegue}
\label{sec:estado_arte_kubernetes}

Los servicios de la plataforma pueden desplegarse de tres formas:

\begin{itemize}
  \item \textbf{Sin orquestador}: cada servicio como proceso o contenedor
  suelto en la máquina. Es la opción más simple de depurar, pero renuncia
  al modelo declarativo de infraestructura versionada y no ejercita
  ninguna competencia de computación en la nube.
  \item \textbf{Docker Swarm~\cite{docker_swarm_docs}}: orquestador más
  simple que Kubernetes, pero prácticamente ausente del soporte oficial
  del ecosistema Big Data. Ni Spark, ni Airflow, ni Superset publican un
  camino de despliegue nativo equivalente al que sí publican para
  Kubernetes.
  \item \textbf{Kubernetes~\cite{cncf_kubernetes}}: el estándar de facto,
  con soporte de primera clase en el resto de la pila elegida.
\end{itemize}

La opción adoptada es Kubernetes, y dentro de él, k3s~\cite{k3s_docs}: una
distribución certificada, empaquetada como binario único con huella de
recursos reducida, apta para una sola máquina de desarrollo sin dejar de
ser Kubernetes real. Esa condición hace posible desplegar sobre una única
máquina local durante el desarrollo, en lugar de sobre un clúster real en
la nube. Evita así el coste de mantener infraestructura en la nube
encendida mientras el sistema todavía cambia con frecuencia.

La arquitectura no depende de esa máquina concreta, porque todos los
servicios se despliegan con \textit{charts} de Helm sobre Kubernetes
estándar. La única excepción es el modo de red \texttt{host-gw}, necesario
bajo WSL2 para evitar un problema conocido de encapsulado UDP en el
\textit{backend} por defecto, una decisión de compatibilidad con el
entorno de desarrollo, no una limitación de Kubernetes. El mismo
despliegue es, por tanto, replicable sobre un clúster real en la nube
cambiando la infraestructura subyacente, no la definición de los
servicios. Ese paso queda documentado como trabajo futuro.

\section{Procesamiento distribuido: Apache Spark frente a sus alternativas}
\label{sec:estado_arte_spark}

El procesamiento de este proyecto tiene forma de lotes programados, no de
flujo continuo de eventos, lo que orienta la elección entre los motores
disponibles:

\begin{itemize}
  \item \textbf{Apache Flink~\cite{apache_flink_docs}}: pensado ante todo
  para baja latencia sobre flujos continuos. Añadiría la complejidad
  operativa de su propia topología de coordinación sin que el proyecto
  llegue a explotar esa ventaja.
  \item \textbf{Dask~\cite{dask_docs} y Ray~\cite{ray_docs}}: motores
  nativos de Python, con curva de adopción más suave. Su soporte de
  Iceberg es bastante menos maduro, y su presencia en arquitecturas
  \textit{lakehouse} de referencia es minoritaria.
  \item \textbf{Apache Spark \cite{zaharia_spark_2016}}: descrito por sus
  autores como un motor unificado para lotes, flujos, SQL y aprendizaje
  automático. Cuenta con la integración más madura del ecosistema con
  Iceberg y con un \textit{backend} nativo sobre Kubernetes
  \cite{apache_spark_k8s_docs}.
\end{itemize}

La opción adoptada es Apache Spark, lanzado con \texttt{spark-submit} en
modo Kubernetes \textit{client} desde el propio orquestador, en lugar de
mediante el \textit{Spark Kubernetes Operator}~\cite{spark_kubernetes_operator_docs}.
Este último exige instalar y depurar un controlador adicional, mientras
que \texttt{spark-submit} en modo \textit{client} sigue programando los
ejecutores como pods reales sin esa pieza extra. Ciertas propiedades de
configuración del controlador de Spark resultaron no operativas en modo
\textit{client} y tuvieron que trasladarse a la especificación del propio
pod, un ajuste documentado en el Capítulo~\ref{cap:infraestructura}.

\section{Orquestación de flujos ETL: Apache Airflow frente a sus alternativas}
\label{sec:estado_arte_airflow}

Cuatro herramientas cubren la orquestación de flujos ETL en el ecosistema
actual:

\begin{itemize}
  \item \textbf{Prefect~\cite{prefect_docs}}: ergonomía de desarrollo
  local moderna, con los flujos definidos como código Python idiomático,
  pero una madurez de despliegue nativo en Kubernetes menor que la de
  Airflow.
  \item \textbf{Dagster~\cite{dagster_docs}}: modelo de activos de datos
  con linaje explícito entre pasos, atractivo para un lago de datos, pero
  no crítico para un \textit{pipeline} de solo tres etapas cuya
  procedencia ya resuelve el propio formato Iceberg.
  \item \textbf{Luigi~\cite{luigi_github}}: simplicidad orientada a
  tareas con dependencias, pero sin interfaz de administración ni un
  ejecutor de Kubernetes nativo comparable.
  \item \textbf{Apache Airflow \cite{apache_airflow_k8s_docs}}: el
  orquestador incumbente, con \textit{chart} de Helm oficial y un
  ejecutor de Kubernetes de primera clase, el \texttt{KubernetesExecutor},
  en el que cada tarea se ejecuta como su propio pod.
\end{itemize}

La opción adoptada es Apache Airflow con su ejecutor de Kubernetes. Frente
a su alternativa madura, el \texttt{CeleryExecutor}, basado en un
intermediario de colas, reutiliza el propio clúster como sustrato de
ejecución en lugar de exigir levantar una cola de mensajes y un almacén
de resultados adicionales.

\section{Capa de consulta interactiva y capa de visualización}
\label{sec:estado_arte_bi}

Trino~\cite{trino_docs}, un motor de consulta SQL federado de baja
latencia sobre tablas Iceberg, sería útil si el proyecto necesitase una
vía de consulta independiente del panel de visualización. No es el caso,
porque todo el acceso a los datos ya se hace desde los propios trabajos de
procesamiento, así que queda como trabajo futuro.

Para la visualización, cuatro opciones son relevantes:

\begin{itemize}
  \item \textbf{Grafana~\cite{grafana_docs}}: la referencia para paneles
  de observabilidad y series temporales, pero encaja peor con gráficos
  analíticos agregados sobre un esquema relacional, la forma real de los
  indicadores de este proyecto.
  \item \textbf{Power BI~\cite{microsoft_power_bi_docs} y
  Tableau~\cite{tableau_docs}}: las herramientas de BI propietarias de
  referencia del mercado, incompatibles con la naturaleza autoalojada y
  de coste cero de la plataforma.
  \item \textbf{Metabase~\cite{metabase_docs}}: alternativa de código
  abierto comparable a Superset, de entrada más simple, pero con un
  modelo de panel menos expresivo para versionar el panel completo como
  código.
  \item \textbf{Apache Superset~\cite{apache_superset_docs}}: permite
  exportar y reconstruir el panel completo desde el repositorio, en lugar
  de depender de un artefacto manual imposible de versionar.
\end{itemize}

La opción adoptada es Apache Superset, conectado a una base de datos
relacional y no a Iceberg. Con Trino ya descartado, materializar las
tablas de \textit{gold}, ya agregadas y pequeñas, en una base de datos
relacional como último paso del \textit{pipeline} da a Superset un
conector nativo sin necesitar ningún servicio intermedio. El
\textit{chart} de Helm oficial de Superset resultó estar obsoleto frente
a la vía de despliegue que el proyecto promueve actualmente. Se gestionó
fijando su versión y construyendo una imagen propia, en lugar de adoptar
un operador todavía en maduración, una decisión que se retoma en el
Capítulo~\ref{cap:visualizacion}.

\Needspace{9cm}
\section{Síntesis y arquitectura resultante}
\label{sec:estado_arte_sintesis}

La Tabla~\ref{tab:tfm_sintesis_decisiones} reúne, para cada capa, la
decisión adoptada, la alternativa principal descartada y el motor de su
justificación.

\begin{table}[H]
  \centering
  \small
  \renewcommand{\arraystretch}{1.2}
  \begin{tabular}{p{0.24\textwidth} p{0.28\textwidth} p{0.4\textwidth}}
    \hline
    \textbf{Decisión} & \textbf{Alternativa principal descartada} & \textbf{Motor de la justificación} \\
    \hline
    Arquitectura \textit{lakehouse} & Almacén y lago separados & Un único sistema para ambas exigencias, sin duplicar operación \\
    Apache Iceberg & Delta Lake, Apache Hudi & Sin dependencia de proveedor, consultable desde otros motores sin cambios \\
    Catálogo Hadoop & Hive Metastore, REST, Nessie & Un único escritor por capa, sin necesidad de versionado de catálogo \\
    MinIO & HDFS, nube pública real & Mismo protocolo que un proveedor real, sin coste ni dependencia externa \\
    Kubernetes k3s, nodo local & Docker Swarm, clúster real en la nube & Soporte de primera clase en el resto de la pila, coste y reproducibilidad durante el desarrollo \\
    Apache Spark, \texttt{spark-submit} sin operador & Apache Flink, Dask, Ray, \textit{Spark Kubernetes Operator} & Patrón por lotes del proyecto, integración madura con Iceberg, menor superficie operativa \\
    Apache Airflow, ejecutor de Kubernetes & Prefect, Dagster, Luigi, \texttt{CeleryExecutor} & Despliegue Kubernetes-nativo maduro, sin servicios adicionales \\
    Superset sobre base de datos relacional, sin Trino & Trino, Grafana, Metabase, BI propietario & Panel como código, sin servicio intermedio adicional \\
    \hline
  \end{tabular}
  \caption{Síntesis de las decisiones arquitectónicas de la plataforma.}
  \label{tab:tfm_sintesis_decisiones}
\end{table}

En conjunto, estas decisiones configuran una arquitectura \textit{lakehouse}
coherente, autoalojada y reproducible desde cero, con soporte de primera
clase entre todos sus componentes y con cada alternativa descartada
documentada como trabajo futuro explícito. El
Capítulo~\ref{cap:infraestructura} describe cómo se despliega
efectivamente. Los Capítulos~\ref{cap:ingesta},~\ref{cap:procesamiento}
y~\ref{cap:visualizacion} retoman, cada uno en su lugar, las decisiones de
estrategia de datos, resolución de identidad y observabilidad propias del
desarrollo de cada capa.
